\subsection*{The exact weighted cost, not a histogram peak (NS-32/33)}

The review of the v1.48 candidate (PR~14, head \texttt{9b2b900})
retains its weighted-cost implication but rejects the claimed cofinal
biconditional and the use of a histogram peak as a bounded continuous
level density. The candidate is not merged by this review. The audit is
archived as \path{audits/2026-09-21-v1.48-pr14-v2.html}.
This subsection gives a self-contained corrected formulation and answers
the zero-statement question. None of the diagnostic optimizations is rerun.

Fix $a>0$, $D=D_a>0$, a prescribed integer head size $N\ge1$, and $X=2\pi N/L$,
$L=2a$, $x=e^L$. Let $\mathcal Z_a$ be the \emph{complete corrected zero field}
of \eqref{eq:v147-corrected-zero-field}; thus $\mathcal Z_a=\beta_a$,
with all nontrivial zeros, the strict-cutoff endpoint and the trivial
terms retained. Its symmetric-height Perron convention and continuous
real limits are unchanged. Define finite positive Borel measures on
$[-D,0]$ by
\begin{equation}\label{eq:v149-measures}
 \widetilde m_a(B)=\frac1X\int_0^X
               1_{\{\mathcal Z_a(\xi)\in B\}}\,d\xi,
 \qquad
 \mu_g(B)=\int_{\{\beta_a\in B\}}|\mathcal F_ag(\xi)|^2\,d\xi.
\end{equation}
For a unit vector both measures have total mass at most one. The negative
fraction and the exact density at almost every regular negative level are
\begin{align}
 \varphi_-(a)&=\frac1X\int_0^X1_{\{\mathcal Z_a(\xi)<0\}}\,d\xi,
 \notag\\
 \widetilde\rho_a(t)&=\frac1X
  \int_{\{\xi\in(0,X):\mathcal Z_a(\xi)=t\}}
       \frac{d\mathcal H^0(\xi)}{|\mathcal Z_a'(\xi)|},
 &R_a&=\mathop{\rm ess\,sup}_{-D<t<0}\widetilde\rho_a(t).
 \label{eq:v149-zero-density}
\end{align}
Here $\mathcal H^0$ counts preimages, so the integral is a sum, and
$R_a$ is allowed to be infinite. The measures have no atoms at levels
$-D$ or $0$. Equation~\eqref{eq:v147-coarea} proves absolute continuity,
not an $L^\infty$ bound for the density.

\begin{proposition}[Critical levels make the peak-density bound vacuous]
\label{prop:v149-critical-density}
Suppose $\xi_0\in(0,X)$ is a critical point with
$t_0=\beta_a(\xi_0)<0$. Then $R_a=\infty$. More precisely, if
\[
 \beta_a(\xi)=t_0+A(\xi-\xi_0)^r+O((\xi-\xi_0)^{r+1}),
 \qquad A\ne0,\quad r\ge2,
\]
one local inverse branch contributes
\begin{equation}\label{eq:v149-critical-asymptotic}
 \frac{1+o(1)}{Xr|A|^{1/r}}
               |t-t_0|^{1/r-1}
\end{equation}
to $\widetilde\rho_a(t)$ on the side reached by that branch.
In particular every complete negative component contained inside
$(0,X)$ supplies an interior minimum and makes $R_a$ infinite.
The same conclusion holds for a one-sided boundary critical point if
its interior branch has negative levels.
\end{proposition}
\begin{proof}
The actual symbol is nonconstant real analytic. Thus the first nonzero
Taylor derivative at a critical point has finite order $r\ge2$.
Inverting on either available monotone branch gives
$|\xi-\xi_0|=(|t-t_0|/|A|)^{1/r}(1+o(1))$ and hence
\eqref{eq:v149-critical-asymptotic}. The exponent is strictly between
$-1$ and $0$: the singularity is integrable but unbounded on sets of
positive level measure. Other preimages add nonnegative terms, so
there is no cancellation in the density. A complete negative component
has zero boundary values and a negative interior minimum, proving the
last assertion, with the boundary case treated one-sidedly.
\end{proof}

Thus there need not be a finite ``most common level'' near zero. The
largest \emph{bin average} in the NS-28 histogram is a different,
resolution-dependent statistic. Near the critical value above its
average on an adjacent bin of width $h$ grows like $h^{1/r-1}$ as
$h\downarrow0$. A finite bin peak cannot upper-bound $R_a$. The theorem
is conditional only on the stated geometric critical point, not on RH;
it does not certify the location of any diagnostic trough.

A finite density upper bound would itself be an all-band statement:
\[
 \widetilde m_a(B)\le R_a|B|\quad\hbox{for every Borel }B\subset[-D,0].
\]
To make $\varphi_-^2/R_a\to\infty$ using a finite upper estimate $r_a$
would suffice to prove, eventually,
\begin{equation}\label{eq:v149-slope-demand}
 \frac1{X_a}\sum_{\mathcal Z_a(\xi)=t}\frac1{|\mathcal Z_a'(\xi)|}
 \le r_a=o(\varphi_-(a)^2)
 \quad\hbox{for almost every }-D_a<t<0.
\end{equation}
A bound on the number of crossings by $n_a$, together with
$|\mathcal Z_a'|\ge s_a>0$ at all their negative regular preimages,
would give $r_a=n_a/(X_as_a)$. A lower slope bound at one level or
one preimage is insufficient; multiplicity of crossings matters too.
Near a negative critical point the displayed uniform slope mechanism
fails. The derivative means the derivative of the \emph{whole} corrected
field, not an unjustified termwise derivative of a conditionally
convergent zero prefix. It can equivalently be evaluated from the exact
finite-prime formula:
\[
 \beta_a'(\xi)=-\tfrac12\Im\psi'(5/4+i\xi/2)
   +2\sum_{1<n<x}\frac{\Lambda(n)\log n}{\sqrt n}\sin(\xi\log n)
   -2\int_0^L t e^{t/2}\sin(\xi t)\,dt.
\]
No replacement $|x^{\rho-1/2}|=1$ is used; that replacement for every
zero would require RH.

\paragraph{The arithmetic of the scalar and its relation to ZLD.}
For $0<\varphi_-\le1$ and $0<R_a<\infty$,
\begin{equation}\label{eq:v149-scalar-arithmetic}
 R_a\ge\frac{\varphi_-}{D_a},\qquad
 0\le\frac{\varphi_-^2}{R_a}\le\varphi_-D_a\le D_a,
 \qquad
 \frac{\varphi_-^2}{R_a}\to\infty
 \ \Longleftrightarrow\ \frac{R_a}{\varphi_-^2}\to0.
\end{equation}
In particular divergence requires $R_a\to0$. A positive lower bound
on $R_a$ prevents it; an upper bound on $R_a$ does \emph{not} prevent
it. Uniform level density has $R_a=\varphi_-/D_a$, and gives the
scalar $\varphi_-D_a$, which can diverge even with $\varphi_-\le1$.
With $R_a=\infty$ the packing lower bound is simply zero. Thus bounded
negative fraction alone is not the reason for rejecting the scalar;
its actual unbounded peak at critical levels is the sharper issue.

ZLD is a lower bound on all level bands; a finite $R_a$ is an upper
bound on all bands. Neither replaces the other as a matter of measure
theory. For example, with $U_I$ the uniform probability measure on
$I$, the measure
$\tfrac14U_{[-D,0]}+\tfrac14U_{[-1,0]}$ satisfies full-range ZLD with
$\kappa=1/4$, but has $\varphi_-=1/2$, $R\to1/4$ and bounded scalar.
Conversely, for fixed $0<\theta\le1$ and $p=D^{-2}$, the measure
\[
 \tfrac12\bigl((1-p)U_{[-\theta D/2,0]}+pU_{[-D,0]}\bigr)
\]
has scalar asymptotic to $\theta D/4$, but no uniform ZLD constant
on $[-\theta D,0]$: the lower part has density $p/(2D)$.
These are full-support measure countermodels to the proposed inference,
not distributions asserted for zeta. No implication, equivalence or
logical independence for the actual zero field is proved by them.

\begin{proposition}[The surviving weighted-cost implication]
\label{prop:v149-weighted-cost}
Let $\mathcal L_K$ be the finite subsets of $[-D,0]$ containing $-D$
and having at most $K\ge1$ elements. For $\mathcal L\in\mathcal L_K$
put $p_{\mathcal L}(t)=\max\{\ell\in\mathcal L:\ell\le t\}$ and define
\begin{equation}\label{eq:v149-qc}
 \mathrm{QC}_K(\widetilde m_a)
 =\inf_{\mathcal L\in\mathcal L_K}
       \int_{[-D,0]}(t-p_{\mathcal L}(t))\,d\widetilde m_a(t).
\end{equation}
This at-most-$K$ formulation avoids any need to assume strict spacing
or attainment of a minimum. Suppose a unit even source-admissible
$g\in\mathcal D(q_a)$ satisfies $q_a[g]\le Q$, $Q\ge0$, and
\begin{equation}\label{eq:v149-wlh}
 \mu_g(B)\ge c'\widetilde m_a(B)
 \quad\hbox{for every Borel }B\subset[-D,0],\qquad c'>0.
\end{equation}
Then every step minorant from \eqref{eq:v131-step-minorant} with at most
$K$ distinct values satisfies
\begin{equation}\label{eq:v149-weighted-cost-bound}
 \eta(m)\ge[c'\mathrm{QC}_K(\widetilde m_a)-Q]_+.
\end{equation}
No bounded density of $\widetilde m_a$ is needed. If it additionally has
a density bounded by a finite positive $R$, then
$\mathrm{QC}_K\ge\varphi_-^2/(2KR)$. For uniform density on $[-D,0]$,
$\mathrm{QC}_K=\varphi_-D/(2K)$, recovering v1.46 with $c=c'\varphi_-$.
\end{proposition}
\begin{proof}
The negative values of $m$ include $-D$ and form a set
$\mathcal L\in\mathcal L_K$. At a negative symbol value $t$ one has
$m\le p_{\mathcal L}(t)$; all discarded loss is nonnegative. Since
$\mu_g-c'\widetilde m_a$ is a finite positive measure, its integral
against the nonnegative Borel function $t-p_{\mathcal L}(t)$ is
nonnegative, by approximation by simple functions. Hence the total
packet loss is at least $c'\mathrm{QC}_K$. Subtract $q_a[g]\le Q$ as
in \eqref{eq:v146-gap-packet}. The infimum is nonincreasing in $K$
because $\mathcal L_K\subset\mathcal L_{K+1}$, so fewer levels cause
no gap in this argument.

Under the additional density bound, write $m_j$ for the mass in each
gap. Packing it nearest the gap's left endpoint at density at most
$R$ costs at least $m_j^2/(2R)$. Summing at most $K$ gaps and using
Cauchy--Schwarz gives $\varphi_-^2/(2KR)$. For uniform density, equal
gaps attain this value, proving the constant exactly.
\end{proof}

\paragraph{The one-level cost is an exact zero functional.}
Put $\mathcal Z_a^-=\max(-\mathcal Z_a,0)$, and keep the global
$D_a$ even when the head range does not attain it. Then
\begin{align}
 A_a:=\mathrm{QC}_1(\widetilde m_a)
 &=\frac1{X_a}\int_{\{0\le\xi\le X_a:\mathcal Z_a(\xi)<0\}}
                         (D_a+\mathcal Z_a(\xi))\,d\xi \notag\\
 &=\varphi_-(a)D_a-\frac1{X_a}\int_0^{X_a}\mathcal Z_a^-(\xi)\,d\xi
 =\int_0^{D_a}\left(\varphi_-(a)-\frac{F_{a,X_a}(t)}{X_a}\right)dt.
 \label{eq:v149-qc1-zero}
\end{align}
This follows directly from $p_{\{-D\}}=-D$ and the layer-cake identity
for the negative part. Thus $A_a\to\infty$ asks that the total
frequency-averaged distance of negative values from the deepest
value diverge. For example it suffices that
$\varphi_-D_a\to\infty$ and
$X_a^{-1}\int\mathcal Z_a^-\le(1-\delta)\varphi_-D_a$ for some fixed
$\delta>0$. It is a statement about values of the corrected zero sum,
not about the number of zeros or a prescribed zero spacing.

There is no conclusion here that $\varphi_-$ tends to a positive limit,
or to zero. Corollary~\ref{cor:v137-scalar-no-go} proves global depth
\emph{only}; it does not prove nested negative sets, increasing measure,
or any fraction for the chosen range. The four observed fractions do
not establish an asymptotic. The choice of $N(a)$ matters: at each fixed
window the negative set is bounded, so choosing an excessively large
prescribed head range can make its fraction arbitrarily small. No
uniform assertion over all choices of $N(a)$ is possible. These points
are independent of RH and do not themselves disprove a positive-fraction
statement for a suitably specified family.

\begin{proposition}[One-level growth does not imply fixed-level growth]
\label{prop:v149-qc1-insufficient}
There is no universal positive constant $\kappa$ such that
$\mathrm{QC}_K(\nu)\ge\kappa\mathrm{QC}_1(\nu)/K$ for all finite
absolutely continuous level measures $\nu$. In fact there are measures
on $[-D,0]$ of mass $1/2$, with support the whole interval, such that
$\mathrm{QC}_1\sim D/2$ but $\mathrm{QC}_2\to0$ as $D\to\infty$.
\end{proposition}
\begin{proof}
Take $D>3$, $h=D^{-1}$, $p=D^{-2}$ and $r=D^{-3}$, and put
\[
 \nu_D=\tfrac12\bigl((1-p-r)U_{[-1-h,-1]}
             +pU_{[-D,-D+h]}+rU_{[-D,0]}\bigr).
\]
The one-level cost is asymptotic to $D/2$, since almost all the mass
sits near $-1$ while the compulsory level is $-D$. The admissible two
levels $\{-D,-1-h\}$ give cost at most $h/4+rD/2$, hence
\[
 \mathrm{QC}_2(\nu_D)\le\frac1{4D}+\frac1{2D^2}\longrightarrow0.
\]
The bridge of mass $r/2$ makes the support all of $[-D,0]$. These are
ordinary absolutely continuous measures. They are not claimed to be
arithmetic level distributions; they disprove the proposed general
quantization inference.
\end{proof}

\begin{assumption}[Exact level quantization growth, QG]
\label{ass:v149-qg}
Fix a cofinal family and a prescribed rule $N(a)$, hence $X_a$.
For every fixed integer $K\ge1$, require
\begin{equation}\label{eq:v149-qg}
 \mathrm{QC}_K(\widetilde m_a)
 =\inf_{\mathcal L\in\mathcal L_K}\frac1{X_a}
   \int_{\{\mathcal Z_a<0\}\cap[0,X_a]}
       (\mathcal Z_a(\xi)-p_{\mathcal L}(\mathcal Z_a(\xi)))\,d\xi
 \longrightarrow\infty.
\end{equation}
Separately require source-admissible unit packets $g_a$ satisfying
\eqref{eq:v149-wlh} with $c'\ge c_0>0$ and $q_a[g_a]\le Q_*<\infty$, $Q_*\ge0$.
Both are named open inputs. Neither follows here from the four diagnostic
windows or from v1.37.
\end{assumption}
Under these hypotheses, a minorant family with bounded level count
$K(a)\le K_0$ has
$\eta(m_a)\ge[c_0\mathrm{QC}_{K_0}(\widetilde m_a)-Q_*]_+\to\infty$.
Thus a uniformly bounded loss forces $K(a)\to\infty$ (apply the same
argument to any subsequence with bounded $K$). This is a sufficient
obstruction, not a necessary condition for all possible proofs of one.
QG says that the corrected zero field's negative values cannot be
approximated from below by any fixed number of levels with bounded
mean error. It needs no bounded peak density.

A stronger quantitative alternative is \emph{quantization shape control}:
for some $\kappa>0$ independent of $a,K$, require
\begin{equation}\label{eq:v149-shape}
 A_a\to\infty,\qquad
 \mathrm{QC}_K(\widetilde m_a)\ge\frac{\kappa A_a}{K}
 \quad\hbox{for all }K\ge1\hbox{ on the cofinal family}.
\end{equation}
With the packet input this gives
$K\ge c_0\kappa A_a/(C+Q_*)$ for loss $\eta\le C$ and $C+Q_*>0$.
If $C+Q_*=0$, positive exact cost instead gives incompatibility.
Linear growth in $D_a$ additionally needs $A_a\ge d_0D_a$ for a fixed
$d_0>0$. Equation~\eqref{eq:v149-qc1-zero} by itself gives none of the
shape inequality: concentration in a few narrow level bands is precisely
the missing obstruction. The empirical relation
$\mathrm{QC}_K\approx\mathrm{QC}_1/(1.8K)$ is not that theorem.

\paragraph{Diagnostic and source constraints after review.}
The archived NS-28 code constrains 20 level bins for raw $N=256$ even
vectors, scans to $\xi=800$ independently of $2\pi N/L$, uses a sampled
depth, and computes a 200-bin surrogate cost with reconstruction levels
on bin edges. The reported finite peak is a peak of bin averages, and
only $K\le32$ costs are evaluated. At $\lambda=4,c'=1$, the saved
returned vector has minimum bin ratio $0.636<1$: its energy is not a
feasible upper bound even for the binned optimization. The dual number
alone supplies no such upper bound. Claims of exact WLH, cofinal
$Q\to0$, and measured costs at hundreds or thousands of levels are
therefore withdrawn as interpretations of those diagnostics; their
raw files are preserved.

For a raw unit vector $f$, projection to
$g=P_af/\sqrt{1-r^2}$ with $r=|\langle f,u_a\rangle|<1$ obeys
\eqref{eq:v147-projected-density}. If $E=q_a[f]$ and
$\delta=(2r+r^2)\|\mathsf W_{e^a}u_a\|$, the form expansion gives
\[
 \frac{E-\delta}{1-r^2}\le q_a[g]\le\frac{E+\delta}{1-r^2}.
\]
For $E>0$, the explicit sufficient checks $r^2\le1/4$ and
$\delta\le E/2$ would keep this within $[E/2,2E]$. No such overlap
check is archived for the LP minimizer; a small source residual alone
does not control the denominator or preserve the weighted level
constraints. The missing $\lambda=3$ computation is to recover that
minimizer in the centered source basis, enclose its overlap and source
tail, its complete energy and the source cross term, and recheck the
projected constraints. It is named, not run. Passing the finite bin
constraints would still not prove the all-Borel hypothesis.

\paragraph{Finding (NS-33).}
The peak-density scalar is the wrong general bound for an exact level
measure containing a negative critical point: it is vacuous there.
The valid alternative is the exact cost, reduced to the named open
input QG plus the admissible bounded-energy weighted packet input.
Replacing it by $\mathrm{QC}_1/K$ requires the additional shape
hypothesis \eqref{eq:v149-shape}; divergence of
$\varphi_-D_a-X_a^{-1}\int\beta_a^-$ alone is insufficient. No cofinal
input is proved unconditionally. All new local identities, implications
and countermodels require no RH. Odd packets satisfy the same measure
argument because their Fourier modulus is even, without the even-source
constraint. This is a limitation of the proposed bound and diagnostic
interpretation, not a failure of the physical form or the valid weighted
concentration criterion. No new window, tail metric, optimization or zero
enclosure is computed. The route remains blocked. G2 and RH remain open.

